%!TEX root = ../LLMSeminar.tex
% ──────────────────────────────────────────────────────────────
%  Section 7 — The Agent Loop
% ──────────────────────────────────────────────────────────────

\section{The agent loop}\label{sec:agent-loop}

This is the technical heart of the seminar.  The thesis is the same
one that appeared in compressed form in \S\ref{sec:agents}, but now
we will state it carefully, write it down, examine it line by line,
and convince ourselves that there is genuinely nothing missing.

\subsection{The whole architecture, in five lines}

\begin{keyresult}[Canonical agent loop]\label{kr:agent-loop}
  \begin{lstlisting}
while not done:
    response   = llm(system_prompt + history)
    tool_calls = parse(response)
    results    = execute(tool_calls)
    history   += [response, results]
  \end{lstlisting}
\end{keyresult}

\begin{remark}
  Every coding agent in production---Claude Code, Cursor, Cline,
  Aider, the recently leaked source of Claude Code itself---is some
  elaboration of this loop.  More tools, better prompts, sharper
  parsing, smarter context management; but no extra ingredient at
  the top of the recipe.  If you can write the five lines, you can
  describe what every commercial coding agent is doing.  This
  alone, on the latest software-engineering surveys, places you
  in the top decile of working developers in 2026.
\end{remark}

\subsection{Reading the loop}

Let us go through it slowly.

\paragraph{Line 1: \texttt{while not done}.}
The agent does not, in any non-trivial sense, decide on its own
when to stop.  It runs until you stop it (Ctrl-C), or until the
harness imposes a limit (a maximum number of turns, a budget cap,
a time-out).  ``Done-ness'' is a property of the harness, not of
the LLM.

\paragraph{Line 2: \texttt{response = llm(...)}.}
This is the only point at which any cognition happens.  The
harness assembles a single string by concatenating the system
prompt (which describes the agent's role and the tools available)
with the running history, ships it to the LLM via an HTTP request,
and receives a string in return.  The function call is stateless;
all relevant context is packed into its input.  Inside the LLM is
a forward pass through a frozen Transformer~\cite{vaswani2017attention}.

\paragraph{Line 3: \texttt{tool\_calls = parse(response)}.}
The response is, mostly, English text, but it can also include
\emph{structured} sub-strings that follow the agreed tool-call
schema.  The harness's job is to pull these out.  The parsing
is an entirely local, non-AI computation: a finite-state matcher
or a JSON parser.  Crucially, it is \emph{not} the LLM that runs
this step---it is plain code on the client side.

\paragraph{Line 4: \texttt{results = execute(tool\_calls)}.}
For each requested tool call, the harness performs the action
and collects a result.  Reading a file returns its contents;
running a shell command returns stdout, stderr, and an exit
code; querying a database returns rows.  None of this is AI.
This step is just system code.

\paragraph{Line 5: \texttt{history += [response, results]}.}
The new response and the tool results are concatenated onto
the running history, so that on the next iteration the LLM sees
both what it said and what came back.  Context grows
monotonically until either the task is done, the loop is
killed, or the context window is exhausted---in which case the
client must compress, summarise, or drop earlier turns
(this is one of the failure modes we discussed in
Seminar~I).

\begin{audience}[Audience question]
  \emph{How does the loop actually terminate?}

  In practice, it doesn't---you give it a quota of turns
  (e.g.\ 25), or you provide a tool such as \texttt{stop()} or
  \texttt{submit\_answer()} that the LLM is encouraged to call
  when satisfied.  Many production agents simply run until the
  user interrupts.  The freedom is yours.
\end{audience}

\subsection{The flow, drawn}

\begin{center}
\begin{tikzpicture}[node distance=1.4cm, every node/.style={align=center}]
  \node[sysbox, fill=spacecadet, text=white, minimum width=4.4cm,
        minimum height=0.95cm, font=\small]
    (sys) {System prompt $+$ history};

  \node[sysbox fill, below=1.1cm of sys, minimum width=4.4cm,
        minimum height=0.95cm, font=\small]
    (llm) {\textbf{LLM}\; {\scriptsize $f : \Str \to \Str$}};

  \node[sysbox, right=2.6cm of llm, minimum height=0.95cm, font=\small]
    (parse) {Parse tool calls};

  \node[sysbox, below=1.1cm of parse, minimum height=0.95cm, font=\small]
    (exec) {Execute tools};

  \draw[sysarrow] (sys) -- (llm);
  \draw[sysarrow] (llm) -- (parse)
    node[midway, above, font=\scriptsize] {response};
  \draw[sysarrow] (parse) -- (exec)
    node[midway, right, font=\scriptsize] {call};
  \draw[sysarrow] (exec.south) -- ++(0,-0.5) -| ([xshift=-0.4cm]llm.south)
    node[near start, below, font=\scriptsize] {results};
  \draw[sysarrow, dashed, color=munsell]
    ([xshift=0.4cm]llm.north) -- ([xshift=0.4cm]sys.south)
    node[midway, right, font=\scriptsize, text=munsell] {append};
\end{tikzpicture}
\end{center}

The dashed return arrow is the only ``magic'': by pasting the
last response and the tool outputs back into the next prompt,
the harness manufactures an illusion of memory and progress.
Internally the function is still memoryless; the recurrence
is provided by the client.

\subsection{The tool-call protocol}

A natural question is how the LLM communicates with the harness.
The mechanism is purely textual.  The vendors---OpenAI,
Anthropic, Google---have largely agreed on a shared schema, the
gist of which is the following.

\paragraph{(a) Tool declaration in the system prompt.}
The harness tells the LLM which tools exist, what each one
does, and what arguments it expects.  A typical declaration:

\begin{lstlisting}
{
  "name": "read_file",
  "description": "Read a UTF-8 text file.",
  "input_schema": {
    "type": "object",
    "properties": {
      "path": { "type": "string" }
    },
    "required": ["path"]
  }
}
\end{lstlisting}

\paragraph{(b) Structured tool request in the response.}
When the model wishes to use a tool, it emits a fragment of
JSON in a stylised form:

\begin{lstlisting}
{
  "type": "tool_use",
  "name": "read_file",
  "input": { "path": "src/main.py" }
}
\end{lstlisting}

The harness extracts this, calls \texttt{read\_file("src/main.py")}
locally, and reports back to the LLM:

\begin{lstlisting}
{
  "type": "tool_result",
  "content": "<file contents here>"
}
\end{lstlisting}

\begin{intuition}[The agreed language]
  Tool use is a small protocol layered on top of natural-language
  output.  It is nothing more than: ``if your output contains a
  fragment of this shape, the harness promises to act on it.''
  The vendors have spent considerable training effort teaching
  models to emit this fragment reliably and at the right moments.
\end{intuition}

\begin{remark}[On reliability]
  These structured emissions used to be unreliable---the
  models would forget brackets, mis-quote keys, dribble JSON
  into prose.  Modern models produce well-formed JSON with
  high reliability.  Errors still occur, because the function
  is nondeterministic, but they have become rare.  Robust
  parsers (with retries and small repairs) handle the residual
  noise.
\end{remark}

\subsection{The minimal toolkit}

A genuinely surprising fact: \emph{two} tools suffice.

\begin{keyresult}[Read and write are sufficient]
  Given only \texttt{read\_file(path)} and \texttt{write\_file(path,
  content)}, an agent can be made to do an extraordinary range of
  useful work---including modifying the harness it itself runs
  inside (\S\ref{sec:self-improvement}).
\end{keyresult}

The argument is the Unix one.  Anything you can compute, you
can compute by reading and writing files; the file system is a
universal substrate for state.  In particular, if the harness's
own source is a file, the agent can read it, write a modified
copy, and---on the next invocation---execute the modification.
We will return to this in the next section.

\begin{audience}[Audience question]
  \emph{Don't you also need an ``execute'' tool?}

  Strictly, no.  But in practice the first thing an agent given
  only \emph{read} and \emph{write} will do is rewrite its own
  source code so that an \emph{execute} tool exists.  ``You only
  need read and write,'' but only as a starting point: the agent
  promotes itself to bash access in short order.
\end{audience}

\subsection{Walking through a real implementation}

A canonical reference implementation in roughly $100$ lines of
Python is given in
Appendix~\ref{sec:appendix-agent}; we walk through the pieces
in lecture.  Key features one sees there:

\begin{itemize}
  \item A \emph{tools array}: a list of JSON objects, one per
    tool, with name, description, and input schema.
  \item A \emph{system prompt}: a short paragraph telling the
    agent what it is, what tools it has, and the conventions it
    should follow (e.g.\ \emph{``preserve the loop and tool
    plumbing in this script''}).
  \item A \emph{tool dispatcher}: a function \texttt{run\_tool(name,
    args)} that looks up the named tool and executes it,
    returning either the result or an error.
  \item A \emph{message log}: a list of objects with \texttt{role}
    (\texttt{"user"}, \texttt{"assistant"}, \texttt{"tool"}) and
    \texttt{content}, which is exactly the conversational history
    that gets re-fed into the function on every turn.
  \item A \emph{loop}: ``up to $N$ turns, call the LLM, parse and
    execute tool calls, append, repeat.''
\end{itemize}

The roles in the message log---\texttt{user},
\texttt{assistant}, \texttt{tool}---matter.  They are not just
cosmetic.

\begin{intuition}[Why the role labels are not optional]
  Recall that the LLM is predicting the next token given the
  context.  If a long history is fed in without indicating who
  said what, the model has to \emph{guess} from contextual cues
  whether each utterance was its own or the user's.  It is good
  at guessing, but the labels save a great deal of energy and
  reduce the chance of confusion.  The vendors have trained
  their models specifically on the \texttt{role}/\texttt{content}
  format; deviating from it leads to garbage in, garbage out.
\end{intuition}

\subsection{Tutorials and references}

\begin{greybox}[Suggested reading]
  \begin{itemize}
    \item Thorsten Ball, \emph{How to Build an Agent}
      (\url{https://ampcode.com})~\cite{ball2025agent}.  Probably
      the most influential tutorial in this genre; the article
      that triggered the cohort of agents we have today.
    \item Anthropic, \emph{Tool use overview}~\cite{anthropictools}
      and \emph{Model Context Protocol}~\cite{mcp2024}.
    \item OpenAI, \emph{Function calling guide}~\cite{openaitools}.
    \item Andrej Karpathy, \emph{nanoGPT}~\cite{karpathynano},
      and Jay Alammar's \emph{Illustrated Transformer}.
  \end{itemize}
\end{greybox}

\begin{historical}[The Claude Code source leak]
  In early 2026 the source code of Anthropic's Claude Code agent
  was inadvertently leaked.  To the disappointment of those
  expecting alchemy, what was found was, broadly, exactly the
  loop above, with a careful tool inventory and a substantial
  but unmagical system prompt.  The reaction in the community
  was: of course.  No-one writes a cleverer loop because there
  is no cleverer loop to write.
\end{historical}

\subsection{Bootstrapping yourself in}

You do not need to write the agent from scratch.  The
canonical recipe, which I recommend, is:

\begin{enumerate}[(1)]
  \item Open a chat client (any modern model will do).
  \item Ask: \emph{``give me a basic Python script for a coding
    agent with read-file and write-file tools, that can improve
    itself.''}
  \item Paste the resulting script into a file, install the
    vendor's SDK, set an API key.
  \item Run it.  First instruction: \emph{``read your own source
    code and improve it.''}
  \item Watch.
\end{enumerate}

The cost is real but modest---five to twenty euros of API
credits to develop a full sense of the dynamics.

\begin{remark}[On API access]
  Claude Code with a personal subscription is convenient but
  closed: the subscription does not provide raw API access, so
  you cannot run \emph{your own} code against it; for that you
  need an API key, which is a separate billing item.  This is
  a deliberate product decision and it is worth understanding
  the difference.  Google and OpenAI offer modest free tiers via
  academic trials, sufficient for the experiments suggested
  here.
\end{remark}
